% Revision 1 evaluation draft. This file deliberately does not overwrite sections/evaluation_final.tex.
% Evidence status is made explicit; no new measurements are claimed by this draft.
\section{Experimental Evaluation}
\label{sec:ev-revision}

\subsection{Evaluation scope and evidence status}

The evaluation distinguishes registry measurements, retrieval measurements, machine-labelled assessments,
simulated verifier/attacker outcomes, synthetic scalability experiments and formal results. The distinction is
material: the public registries do not provide human-validated effect labels or differentiated quality-of-service
(QoS) measurements, and the current study does not execute the selected toolkits against live MCP servers.

The merged snapshot contains 76 servers and 569 tools. Only 44 tools publish MCP annotations, all declaring
\texttt{readOnlyHint=true}. The two reference label sets are LLM-generated: \(J_A\) contains 131 tools labelled
state-modifying, while \(J_B\) contains 120. Importantly, the unsafe set of \(J_B\) is contained in that of \(J_A\).
The simulated verifier uses \(J_A\), whereas the scoring step uses \(J_B\). Thus, under the ideal simulated verifier,
zero admission of a \(J_B\)-unsafe tool follows from the label-set construction; it is not independent evidence that
a real verifier has zero false accepts. Neither label set is human ground truth.

\subsection{Retrieval and composition}

On the reported 569-tool E8 registry, PaMCP's goal-contract retrieval obtains 23.0\% Hit@1, compared with
50.0\% for MCP-Zero and 38.7\% for raw-text reciprocal-rank fusion. This is a negative result for the current
structural retrieval configuration. The existing diagnosis identifies disagreement between goal output types and
tool output types as a major source of ranking errors. Therefore, the current evidence does not support a claim
that structural ranking improves full-registry retrieval.

For the 94 LiveMCPBench tasks in E9, PaMCP selects smaller toolkits (1.47 tools on average) but obtains group
recall of 0.169. Per-subtask retrieval obtains 0.249, RAG-MCP 0.293, COLT 0.295 and ToolScope 0.361 under the
reported configurations. Non-inferiority to per-subtask retrieval is not established (\(\Delta=-0.080\),
95\% CI \([-0.166,0.002]\)). The result supports a toolkit-size/context trade-off, not superior functional coverage.

E10 reports 67.0\% joint feasibility for PaMCP under its separate evaluator. This is a labelled proxy: the evaluator
uses the analyst's goal records and LLM-generated \(J_B\) labels, and it simplifies semantic type and effect-scope
reasoning. The comparison with eligible-only per-subtask retrieval is not significant after the stated correction
(\(\Delta=0.043,\ p_H=1.0\)). The result does not establish improved end-to-end task success.

\subsection{Trust and policy evaluation}

E4--E7 explore annotation, lexical and simulated verifier encodings. Their findings describe outcomes under the
chosen label-derived verifier and error simulator. Utility can be measured under these conditions, but verifier
false-acceptance rates are simulator parameters rather than measured rates from executed server probes. In
particular, zero exposure under the error-free verifier must not be interpreted independently of
\(U_{J_B}\subseteq U_{J_A}\).

The current E14 experiment is an offline reachable-action analysis. No injected action is executed. Moreover,
the previous poisoning statistic assigned zero exposure to selected sanitised variants rather than scanning the
actual serialized manifest, and its task-level attack-set builder omitted tampering for tasks with any required
effect. These values are not retained as evidence of sanitisation or runtime attack prevention. Source inspection
of the current sanitizer shows that it reconstructs the description but retains provider-controlled tool names
and argument keys in the model-visible manifest. Revision 1 adds a canary audit for these fields, an alias-manifest
prototype with a server-side reverse map, and a scope-aware re-analysis; their generated output must be reviewed
before any revised numeric result is inserted here. Because the revised estimand differs from the original F5,
the re-analysis is exploratory and does not inherit the original confirmatory status.

Task effect requirements are interpreted at their declared scope. For example,
\texttt{write(/reports)} does not authorize \texttt{write(/secrets)} or an unknown, unscoped write. When the
reference labels contain only an effect class and no target scope, the outcome is reported as indeterminate rather
than safe.

\subsection{Profile, scale, evolution and repair}

The model profile experiment uses 64 calibration items across four operation classes and the tested input sizes.
It supports observations about these items and models only; it does not validate the full conditional reliability
function over all transformation classes, languages and difficulty levels.

The 10k/100k scalability experiments use cloned contracts derived from the 569-tool registry. They are synthetic
scaling experiments. The one-million-entry result measures bitmap eligibility only and must not be described as
full-engine one-million-tool scalability. The eager pairwise risk relation also contributes materially to engine
construction cost.

The current registry provides no differentiated prices or measured latency values, so the QoS weights are not
meaningfully exercised. The generated exact-versus-greedy cover experiments establish behaviour only for their
specified generated instances and node budget. E17 evaluates deterministic repair operators; the LLM assistant
repair role was not run. E13 end-to-end execution, E19 role configurations and E21 experience reuse remain
unmeasured.

\subsection{Threats to validity}

\emph{Construct validity.} Tool types, effects and private-source labels are LLM-generated. Agreement between
models is not a substitute for human or runtime ground truth. Joint feasibility and attack reachability are proxies,
not executed task outcomes.

\emph{Internal validity.} The simulated verifier's evidence source is label A, and label B's unsafe set is contained
in A's. The ideal verifier's zero unsafe admission is consequently construction-dependent. The baseline
implementations include documented adaptations; those results should not be conflated with faithful reproductions
of the original systems. The current retrieval comparison also has a negative result for PaMCP on the merged
registry.

\emph{External validity.} The registry snapshot has 569 tools and sparse annotations. Larger registries are
constructed by cloning. QoS data are absent. No live MCP task-execution benchmark or human-annotated tool
ground truth is included.

\subsection{Summary}

The current evidence establishes the feasibility of the reference implementation and provides useful measurements
of registry metadata, retrieval behaviour, toolkit size, labelled feasibility proxies, index maintenance and
synthetic scaling. It does not yet establish end-to-end task success, real-world verifier accuracy, executed
attack prevention, QoS advantage, LLM-assisted repair effectiveness or experience-reuse benefit. Those claims
require the additional experiments specified in the Revision 1 protocol.